% \subsection{Protein to Phage Aggregation}
% - The raw data is at the protein level but conformal prediction operates at the phage level — aggregation is needed\\
% - NaN-aware mean pooling: for each phage and each functional score column, compute the mean across all proteins of that phage that 
%   have a non-NaN value for that column. Proteins with NaN for a given column are excluded from the mean for that column only\\
% - Why not replace NaN with zero before pooling: treating an unannotated protein as having zero evidence would push the phage's 
%   aggregated score artificially toward zero, misrepresenting its biology\\ 
% - After aggregation: result is a matrix of shape (number of phages) × 40, still with NaNs in cells where none of the phage's 
%   proteins were annotated with that function


% \subsubsection{NaN handling after aggregation}
% - After mean pooling, some phages still have NaN in certain columns -> this happens when none of their proteins were annotated with 
%   a particular function\\
% - Strategy: fill remaining NaNs using column-wise medians computed strictly on the training set (never the test set, to prevent 
%   data leakage)
% - The median cap: any training median above 0.6 is capped at 0.5 -> this prevents an imputed value from falsely pushing a phage 
%   toward a particular class
% - any remaining NaNs after median imputation are filled with 0.5 (neutral value)


\subsection{Protein-to-Phage Aggregation and NaN-Aware Pooling}
\label{sec4.2: aggregation}
\noindent The base classifiers only score individual proteins $p \in \mathcal{P}_i$, while host infection is a property of the 
phage as a whole, so we need to aggregate the protein level scores into a single phage-level score vector. For each phage 
$i \in \{1, \ldots, N_{\text{phage}}\}$ ($N_{\text{phage}} = 18{,}474$) with protein set $\mathcal{P}_i$, let $M_{p,k} = 
\mathbf{1}(\text{protein } p \text{ has a valid prediction for column } k)$ denote the protein-level mask. We compute the 
aggregation using running sums and counts per phage, this avoids holding the full $1.85 \times 10^6$-row protein matrix in memory.

\vspace{0.15cm}
\noindent For each of the $K$ functional dimensions, we get the aggregated phage level score using dimension wise, NaN aware 
mean pooling:
\begin{equation}
\label{eq: nan-mean-pooling}
s_{i,k} = \begin{cases}
\dfrac{\sum_{p \in \mathcal{P}_i} s_{p,k} \cdot M_{p,k}}{\sum_{p \in \mathcal{P}_i} M_{p,k}} & \text{if } 
\sum_{p \in \mathcal{P}_i} M_{p,k} > 0, \\[10pt]
\text{NA} & \text{otherwise.}
\end{cases}
\end{equation}

\noindent Equation~\eqref{eq: nan-mean-pooling} excludes missing entries ($\texttt{NA}$) from both the numerator and denominator, 
and takes the average strictly over observed proteins. A score is $\texttt{NA}$ if and only if no protein of phage $i$ has a 
valid prediction for column $k$.


\vspace{0.15cm}
\noindent The phage-level mask follows directly:
\begin{equation}
\label{eq: phage-mask}
m_{i,k} = \mathbf{1}\left( \sum_{p \in \mathcal{P}_i} M_{p,k} > 0 \right).
\end{equation}
i.e. $m_{i,k}=1$ if at least one protein of phage $i$ has an observed score for column $k$, otherwise, it is $0$. This reduces 
the protein data to a phage-level matrix $\mathbf{S} \in ([0,1]\cup\{\text{NA}\})^{N_{\text{phage}} \times K}$ paired with a 
binary mask for missingness $\mathbf{M} \in \{0,1\}^{N_{\text{phage}} \times K}$.

\vspace{0.15cm}
\noindent As we have already established in Section~\ref{sec3: methodology}, our conformal envelope constructions operate 
natively on the paired observation $(\mathbf{s}_i, \mathbf{m}_i)$.  This preserves the structural missingness without 
fabricating values in place of it. Consequently, all conformal classification results reported in Section~\ref{sec5: experiments} 
rely entirely on this native representation, and require no imputation after aggregation. The task dependent rules that determine 
which protein predictions count as valid and thus define $M_{p,k}$ are detailed in Section~\ref{sec4.3: contamination fix}.